\section{A sharper logarithmic bound from derivative variation}
\label{sec:bv-activations}
A signed-hinge decomposition avoids the additional sample-size logarithms in the preceding general Lipschitz argument when the activation derivative has finite total variation. This gives an alternative proof for the standard activations below, including nonzero activation offsets.

Assume $\phi:\R\to\R$ is globally $K$-Lipschitz and its almost-everywhere derivative has a bounded-variation representative $g$ on $\R$. Write
\[
 a=\phi(0),\qquad V_1=\operatorname{TV}(g),\qquad
 M=\max\{1,K+V_1\},\qquad T=2+|a|+2M.
\]
A Lipschitz function is locally absolutely continuous, so its derivative determines all increments by integration. We may take $|g|\le K$, including its one-sided limits. ReLU is allowed: its derivative has one jump and the corresponding distributional second derivative is a unit point mass. The following safe constants illustrate the scope; numerical sharpness is not needed.
\begin{center}
\begin{tabular}{llll}
\toprule
Activation & $\phi(0)$ & Lipschitz bound $K$ & derivative variation bound $V_1$ \\
\midrule
Tanh & $0$ & $1$ & $2$ \\
ReLU & $0$ & $1$ & $1$ \\
GELU, $x\Phi(x)$ & $0$ & $2$ & $3$ \\
Sigmoid, $(1+e^{-x})^{-1}$ & $1/2$ & $1/4$ & $1/2$ \\
Softplus, $\log(1+e^x)$ & $\log2$ & $1$ & $1$ \\
\bottomrule
\end{tabular}
\end{center}
For GELU, $\phi''(x)=(2-x^2)\varphi(x)$, where $\varphi$ is the standard Gaussian density, so $\int|\phi''|\le\E(2+G^2)=3$; also $|\phi'|=|\Phi+x\varphi|\le2$. The other variation bounds follow from monotonicity on each side of the derivative's single maximum, or from a single jump.

\subsection{Normalized activations are mixtures of centered hinges}
For $t\in[-1,1]$, put
\[
 h_t(u)=\operatorname{ReLU}(u-t)-\operatorname{ReLU}(-t),
 \qquad -1\le u\le1.
\]
Each $h_t$ vanishes at zero, is 1-Lipschitz, and is bounded by one. The identity function is $h_{-1}$ on this interval.

\begin{lemma}[A scale-independent signed mixture]
\label{lem:bv-profile}
For every $s>0$, the function
\[
 u\longmapsto\frac{\phi(su)-\phi(0)}{Ms},\qquad u\in[-1,1],
\]
belongs to the closed absolute convex hull of the centered hinges $h_t$, $t\in[-1,1]$.
\end{lemma}
\begin{proof}
Let $\mu=dg$ be the finite signed distributional derivative measure, and let $\nu_s$ be its restriction to $(-s,s)$ pushed forward by $x\mapsto x/s$. The fundamental theorem for bounded-variation functions gives, on $[-1,1]$,
\begin{equation}
 \frac{\phi(su)-\phi(0)}s
   =g(-s+)u+\int_{(-1,1)}h_t(u)\,d\nu_s(t).
 \label{eq:bv-representation}
\end{equation}
One way to check the identity is to differentiate both sides almost everywhere: the derivative on the right is $g(-s+)+\nu_s((-1,u))=g(su)$ except at a null set, and both sides vanish at zero. Endpoint atoms are absorbed into the one-sided slope or do not affect the interval's interior. Continuity extends equality to the endpoints.

The absolute coefficient mass in \eqref{eq:bv-representation} is at most $K+V_1\le M$. Since $u=h_{-1}(u)$, division by $M$ gives the assertion. The integral is approximated by finite signed mixtures uniformly: $t\mapsto h_t$ is continuous in uniform norm and its signed mixing measure is finite. Zero belongs to the absolute convex hull and supplies any unused mass.
\end{proof}

\subsection{A width-independent capacity estimate}
On a list of $n$ inputs, let $\mathcal A_0$ be the closed absolute convex hull of the $d+1$ coordinate sample vectors of $\bar x=(1,x)$. Recursively let
\[
 \mathcal A_\ell=\overline{\operatorname{aconv}}\left(
 \{\mathbf1\}\cup\{h_t(f):f\in\mathcal A_{\ell-1},\ -1\le t\le1\}\right).
\]
Every class is symmetric, contains zero, and lies in $[-1,1]^n$.

\begin{lemma}[Capacity of the normalized feature classes]
\label{lem:bv-feature-capacity}
For $H\ge0$,
\begin{equation}
 \operatorname{Rad}_n(\mathcal A_H)
 \le\frac{2^H}{\sqrt n}
       \left[\sqrt{2\log(2(d+1))}+4\right].
 \label{eq:bv-feature-capacity}
\end{equation}
\end{lemma}
\begin{proof}
Write $\operatorname{Rad}_n(\mathcal B)=n^{-1}\E_\epsilon\sup_{f\in\mathcal B}\sum_i\epsilon_if_i$. For a symmetric class this also equals the version with an absolute value. The starting class is the convex hull of $2(d+1)$ vectors of Euclidean norm at most $\sqrt n$, so its complexity is at most $\sqrt{2\log(2(d+1))/n}$.

Let $\mathcal A$ be any symmetric class as above and $\mathcal H(\mathcal A)=\{h_t(f):f\in\mathcal A,t\in[-1,1]\}$. Scalar Rademacher contraction for ReLU, with constant one for the ordinary supremum, gives
\[
 \operatorname{Rad}_n(\{\operatorname{ReLU}(f-t):f\in\mathcal A,t\in[-1,1]\})
 \le\operatorname{Rad}_n(\mathcal A)+\E|\textstyle\sum_i\epsilon_i|/n.
\]
The subtracted scalar ranges over $[-1,0]$. Its Rademacher complexity is $\E|\sum_i\epsilon_i|/(2n)$. Splitting the supremum therefore gives
\[
 \operatorname{Rad}_n(\mathcal H(\mathcal A))
 \le\operatorname{Rad}_n(\mathcal A)+\frac3{2\sqrt n}.
\]
This class contains zero. Its positive and negative suprema are nonnegative and have equal expectations, so absolute convexification costs at most a factor two. Adding the constant vector costs at most $1/\sqrt n$. Thus
\[
 \operatorname{Rad}_n(\mathcal A_\ell)
 \le2\operatorname{Rad}_n(\mathcal A_{\ell-1})+4/\sqrt n.
\]
Iterating proves \eqref{eq:bv-feature-capacity}. All classes on a fixed finite input list are compact after closure, so no separability issue enters these finite-dimensional suprema.
\end{proof}

Use the usual full weighted energy
\[
 Q_N=\frac dN\|W^1\|_F^2+\sum_{\ell=2}^H\|W^\ell\|_F^2+N\|v\|_2^2,
 \qquad L=H+1.
\]
\begin{theorem}[A common energy-ball capacity for all BV-derivative activations]
\label{thm:bv-capacity}
For every $q\ge0$, let $\mathcal Z_q$ be the union of raw network outputs over all widths with $Q_N\le q$. Then
\begin{align}
 \sup_{z\in\mathcal Z_q}\|z\|_\infty
 &\le T^H(1+q)^{L/2},\notag\\
 \operatorname{Rad}_n(\mathcal Z_q)
 &\le(2T)^H\left[\sqrt{2\log(2(d+1))}+4\right]
       \frac{(1+q)^{L/2}}{\sqrt n}.
 \label{eq:bv-capacity}
\end{align}
There is no dependence on learner width in either bound.
\end{theorem}
\begin{proof}
Define nonnegative coordinate budgets
\[
 \rho^0=\mathbf1_{d+1},\qquad
 \rho^\ell=|a|\mathbf1_N+M|W^\ell|\rho^{\ell-1}.
\]
An induction shows that hidden neuron $j$'s sample vector lies in $\rho_j^\ell\mathcal A_\ell$. Indeed its preactivation is $s f$ with $s=(|W^\ell|\rho^{\ell-1})_j$ and $f\in\mathcal A_{\ell-1}$. If $s>0$, Lemma~\ref{lem:bv-profile} and the constant component $a$ put its activation in $(|a|+Ms)\mathcal A_\ell$. If $s=0$, its activation is exactly $a$, and the same conclusion holds. A zero budget gives a zero neuron. The output sample vector consequently belongs to $P\mathcal A_H$, where $P=|v|^{\mathsf T}\rho^H$.

Put $b_\ell=\|\rho^\ell\|_2/\sqrt N$. The full energy gives
\[
 b_1\le|a|+M\sqrt{(d+1)/d}\sqrt{Q_N},\qquad
 b_\ell\le|a|+M\sqrt{Q_N}\,b_{\ell-1}.
\]
Since $d\ge1$ and $T\ge|a|+2M$, induction yields
$b_\ell\le T^\ell(1+Q_N)^{\ell/2}$. Therefore
\[
 P\le\sqrt N\|v\|_2 b_H
 \le T^H(1+Q_N)^{L/2}.
\]
The amplitude follows from $\mathcal A_H\subset[-1,1]^n$, and the Rademacher estimate from Lemma~\ref{lem:bv-feature-capacity}. The argument is uniform over every input list and every width. Countable rational parameter approximations with strict energy inequalities, followed by continuity, give a countable determining subclass of the actual network union.
\end{proof}

\subsection{Gaussian teacher neighborhoods without a zero-origin assumption}
At canonical Gaussian scale $s=N^{-\gamma}$, the first, internal, and readout variances are $Ns/d,s,s/N$. Assume $Ns\le1$ and $sD_N\le D_0$, where $D_N=N(d+2)+(H-1)N^2$ and $D_0=d+H+1$. These inequalities hold for every $\gamma\ge2$, and for every $\gamma\ge1$ when $H=1$.

Balanced cloning of a width-$m$, coefficient-$B$ teacher works for every activation, including nonzero offsets: every neuron belongs to a nonempty clone group and has exactly its teacher preactivation. It gives
\[
 E=d(d+1)B^2+2Hm^2B^2
\]
as an upper bound on full weighted energy. By Theorem~\ref{thm:bv-capacity}, take $U=T^H(1+E)^{L/2}$ as an output amplitude bound.

\begin{lemma}[A comparison measure for a general Lipschitz activation]
\label{lem:bv-comparison}
Put
\[
 B_0=1+|a|+2\max\{1,K\}(\sqrt E+1),\qquad
 K_g=2L B_0^H,\qquad
 A_{g,n}=E+1+2D_0\log[4K_g(U+1)n].
\]
At every eligible width and scale, and every $\lambda\ge1/n$,
\begin{equation}
 F_S(\lambda)\le\widehat R_S(g)+\lambda A_{g,n}.
 \label{eq:bv-comparison}
\end{equation}
\end{lemma}
\begin{proof}
Perturb the exact clone by a Gaussian with covariance $\varepsilon^2$ times the prior covariance. Write $b_\ell$ for the normalized $L^2$ size of its random hidden state and $\Delta_\ell$ for its normalized $L^2$ difference from the exact state. Independence of the current Gaussian perturbation, the inequality $Ns\le1$, and $|\phi(t)|\le|a|+K|t|$ give
\[
 b_1\le|a|+\sqrt2K(\sqrt E+\varepsilon),\qquad
 b_\ell\le|a|+K(\sqrt E+\varepsilon)b_{\ell-1}.
\]
For $0<\varepsilon\le1$, these imply $b_\ell\le B_0^\ell$. Lipschitzness also gives
\[
 \Delta_1\le\sqrt2K\varepsilon,\qquad
 \Delta_\ell\le K\sqrt E\Delta_{\ell-1}+K\varepsilon b_{\ell-1}
 \le2\max\{1,K\}\ell\varepsilon B_0^{\ell-1}.
\]
The readout contributes at most $\sqrt E\Delta_H+\varepsilon b_H\le\varepsilon K_g$ to output root mean square error. Only $s\le1$ is used for the readout noise term.

Choose $\varepsilon=[4K_g(U+1)n]^{-1}$. The expected empirical loss is at most $\widehat R_S(g)+1/n$. The exact Gaussian relative-entropy formula gives normalized cost at most
\[
 2s\operatorname{KL}(\rho\Vert\mu)
 \le E+2D_0\log(1/\varepsilon).
\]
Substitute this comparison Gaussian into the free-energy variational identity and use $1/n\le\lambda$. This proves \eqref{eq:bv-comparison}.
\end{proof}

\subsection{The resulting joint activation oracle}
The preceding two results verify the full-energy Gaussian free-energy theorem with
\[
 r=L,\qquad V=T^H,\qquad
 \kappa=(2T)^H\left[\sqrt{2\log(2(d+1))}+4\right],
 \qquad D_0=d+H+1.
\]
A uniform prior-risk bound is also available: with $A_0=\max\{1,K,|a|\}$, conditional Gaussian second moments give
$\E_\mu\widehat R_S\le1+2(4A_0^2)^H$.
Thus the same single free-energy root, with these constants, yields
\begin{equation}
 R(\widehat f_n)\le\inf_g\left\{R(g)+\widetilde O\left(
 A_{g,n}R(g)^{1-1/L}n^{-1/L}+A_{g,n}^L/n\right)\right\}.
 \label{eq:bv-oracle}
\end{equation}
The precise simultaneous interpretation uses the population minimizer of each integer width and coefficient cap and a countable Bernstein union, as in the Gaussian theorem. The prediction is the ordinary mean of raw posterior outputs, and the stronger posterior-average risk obeys the same bound. The event is simultaneous over all eligible widths and every $\gamma\ge2$, or every $\gamma\ge1$ at one hidden layer. The learner knows the activation and architectural depth, but not the comparison cap or risk. No first-layer zero-origin assumption, support restriction, or architecture mixture is used.

At fixed dimension and depth, $A_{g,n}=\widetilde O(E)$, so an exactly fitting teacher gives $\widetilde O(E^L/n)$. Dimension enters polynomially through the explicit comparison and capacity constants. The factors exponential in $H$ are displayed; this argument does not assert polynomial dependence jointly on width and depth. Tanh's sharper depth exponent and ReLU's smaller homogeneous implementation budget remain useful refinements of this common argument.
